\section{无穷级数}
\subsection{常数项级数}

\begin{mydef}{常数项级数}{ConstantTermSeries}
    若级数\(
        \sum_{n=1}^{\infty} a_n
    \)中\(a_n\)为常数，则称该级数为常数项级数。
\end{mydef}

\begin{mydef}{部分和,收敛与余项}{PartialSum}
    级数\(
        \sum_{n=1}^{\infty} a_n
    \)的部分和为
    \[
        S_n = a_1 + a_2 + \cdots + a_n
    \]

    若\(S_n\)收敛，则称该级数收敛；若\(S_n\)发散，则称该级数发散。
    若级数\(
        \sum_{n=1}^{\infty} a_n
    \)收敛，则其和为
    \[
        S = \lim_{n \to \infty} S_{n} = \lim_{n \to \infty} a_n
    \]
    其余项为
    \[
        R_n = S - S_n
    \]
    \end{mydef}
\begin{conclusion}{收敛级数的基本性质}{ConvergenceProperties}
    \begin{enumerate}
        \item 收敛级数的必要条件：级数\(
            \sum_{n=1}^{\infty} a_n
        \)收敛，则\(
            \lim\limits_{n \to \infty} a_n = 0
        \)
        \item 若两个级数\(
            \sum_{n=1}^{\infty} a_n
        \)和\(
            \sum_{n=1}^{\infty} b_n
        \)都收敛，则\(
            \sum_{n=1}^{\infty} (\lambda a_n + \mu b_n)
        \)也收敛。\\
        若一个级数收敛，另一个发散，则\(
            \sum_{n=1}^{\infty} (\lambda a_n \pm \mu b_n)
        \)发散。\\
        若两个级数都发散，则
        \(
            \sum_{n=1}^{\infty} (\lambda a_n \pm \mu b_n)
        \)
        敛散性需进行讨论。
        \item 若级数收敛，则其子级数也收敛。
        \item 在级数中添加或减少有限项不影响敛散性
    \end{enumerate}
\end{conclusion}

\begin{formula}{几个重要的级数}{ImportantSeries}
    \begin{enumerate}
        \item 几何级数
        \(
            \sum_{n=1}^{\infty} q^n
        \)
        当且仅当
        \(
            |q| < 1
        \)时收敛，当
        \(
            |q| \geq 1
        \)时发散。
        \item p级数
        \(
            \sum_{n=1}^{\infty} \frac{1}{n^p}
        \)
        当且仅当
        \(
            p > 1
        \)时收敛，当
        \(
            p \leq 1
        \)时发散。
        \item \(
            \sum_{n=2}^{\infty} \frac{1}{n^p \ln^q n}
        \)
        当
        \(
            p > 1\text{或} p =1 \text{且} q > 1
        \)时收敛，当
        \(
            p \leq 1 \text{或} p =1 \text{且} q \leq 1
        \)时发散。
    \end{enumerate}
\end{formula}

    \begin{mydef}{正项级数}{PositiveTermSeries}
        若级数\(
            \sum_{n=1}^{\infty} a_n
        \)中\(a_n\)全为正数，则称该级数为正项级数。
        正项级数的部分和数列单调非减。
        因此正项级数收敛的充要条件是它的部分和数列有界，即
        \[
            \text{有界} \Leftrightarrow \text{收敛}
        \]
    \end{mydef}
    \subsubsection{正项级数的收敛性的判别法}
    \begin{conclusion}{比较判别法}{ComparisonTest}
        设\(cv_n \geq u_n \geq 0(n \geq N), c > 0\)
        若\(
            \sum\limits_{n = 1}^{\infty} v_n
        \)收敛，则\(
            \sum\limits_{n = 1}^{\infty} u_n
        \)也收敛。
        若\(
            \sum\limits_{n = 1}^{\infty} u_n
        \)发散，则\(
            \sum\limits_{n = 1}^{\infty} v_n
        \)也发散。
        比较原理的极限形式为 \\
        设\(u_n \geq 0\), \(v_n \geq 0\), 若\(
            \lim\limits_{n \to \infty} \frac{v_n}{u_n} = A
        \),则有
        \begin{enumerate}
            \item 若 \(0 < A < \infty\), 则\(
                \sum\limits_{n = 1}^{\infty} u_n
            \)与\(
                \sum\limits_{n = 1}^{\infty} v_n
            \)具有相同的敛散性。
            \item 若 \(A = 0\), 则\(v_n\)比\(u_n\)高阶，\(
                v_n < u_n
            \)
            若\(
                \sum\limits_{n = 1}^{\infty} u_n
            \)收敛，则\(
                \sum\limits_{n = 1}^{\infty} v_n
            \)也收敛。
            \item 若 \(A = \infty\), 则\(u_n\)比\(v_n\)高阶，\(
                u_n < v_n
            \)
            若\(
                \sum\limits_{n = 1}^{\infty} u_n
            \)发散，则\(
                \sum\limits_{n = 1}^{\infty} v_n
            \)也发散。
        \end{enumerate}
    \end{conclusion}
    \begin{conclusion}{比值与根值判别法(与几何级数比较)}{RatioAndRootTest}
       \begin{enumerate}
        \item 比值判别法(达朗贝尔判别法)
        \[
            \lim\limits_{n \to \infty} \frac{u_{n+1}}{u_n} = r = 
            \begin{cases}
                r < 1 & \text{收敛} \\
                r > 1 & \text{发散} \\
                = 1 & \text{此判别法失效}
            \end{cases}
        \]
        且 \(
            \lim\limits_{n \to \infty} u_n = + \infty
        \)
        \item 根值判别法
        \[
            \lim\limits_{n \to \infty} \sqrt[n]{u_n} = R = 
            \begin{cases}
                R < 1 & \text{收敛} \\
                R > 1 & \text{发散} \\
                = 1 & \text{此判别法失效}
            \end{cases}
        \]
         且 \(
            \lim\limits_{n \to \infty} u_n = + \infty
        \)
       \end{enumerate}
    \end{conclusion}
    \begin{conclusion}{与\(p\)级数比较，确定无穷小\(u_n\)关于\(\frac{1}{n}\)的阶}
        设\(u_n > 0\),\(
            \lim\limits_{n \to \infty} \frac{u_n}{\frac{1}{n^p}} = A
        \)
        若
        \begin{enumerate}
            \item \(
                0  < A < +\infty, p > 1
            \)即当\(n \to \infty\)时，\(u_n\)关于\(\frac{1}{n}\)的阶为\(p\),是高阶无穷小。
            则\(
                \sum\limits_{n = 1}^{\infty} u_n
            \)收敛。
            \item \(
                0 < A < +\infty, p \leq 1
            \)即当\(n \to \infty\)时，\(u_n\)关于\(\frac{1}{n}\)的低阶无穷小，
            则\(
                \sum\limits_{n = 1}^{\infty} u_n
            \)发散。
        \end{enumerate}
    \end{conclusion}

    \begin{conclusion}{积分判别法}{IntegralTest}
        设\(
            \sum\limits_{n = 1}^{\infty} u_n
        \)为正项级数,若\(
            \exists
        \)单调下降的函数\(u_n  = f(n)\),则级数\(
            \sum\limits_{n = 1}^{\infty} u_n
        \)收敛的充分必要条件是\(
            \int_{1}^{\infty} u_n d n 
        \)收敛
    \end{conclusion}
    \subsubsection{交错级数}
    \begin{conclusion}{交错级数的定义}{InterlacedSeries}
        若\(u_n \geq 0\)，则称\(
            \sum\limits_{n = 1}^{\infty} {(-1)^{n-1} u_n}
        \)为交错级数。
    \end{conclusion}

    \begin{conclusion}{莱布尼茨判别法}{LeibnizTest}
        设\(
            \sum\limits_{n = 1}^{\infty} {(-1)^{n-1} u_n}
        \)为交错级数,满足以下条件
        \begin{enumerate}
            \item \(
                u_n \geq u_{n+1}, n \geq N \geq 1
            \)
            \item \(
                \lim\limits_{n \to \infty} u_n = 0
            \)
        \end{enumerate}
        则\(
            \sum\limits_{n = 1}^{\infty} {(-1)^{n-1} u_n}
        \)收敛。且其和\(0 < \sum\limits_{n = 1}^{\infty} {(-1)^{n-1} u_n} < u_1\)余项
        \(
            |r^n| < u_{n+1}
        \)
    \end{conclusion}
    \begin{mydef}{绝对收敛与条件收敛}{AbsoluteConvergenceAndConditionalConvergence}
        设\(\sum_{n = 1}^{\infty}u_n\)为任意项级数 \\
        若\(
            \sum_{n = 1}^{\infty} |u_n|
        \)收敛,则称\(\sum_{n = 1}^{\infty}u_n\)绝对收敛。 \\
        若\(
            \sum_{n = 1}^{\infty} |u_n|
        \)发散,而\(
            \sum_{n = 1}^{\infty} u_n
        \)收敛,则称\(\sum_{n = 1}^{\infty}u_n\)为条件收敛。
    \end{mydef}
    \begin{conclusion}{关于绝对收敛与条件收敛的结论}{conclusion about AbsoluteConvergenceAndConditionalConvergence}
        \begin{enumerate}
            \item 绝对收敛的级数一定收敛。即若
            \(
                \sum_{n = 1}^{\infty} |u_n|
            \)收敛,则\(
                \sum_{n = 1}^{\infty} u_n
            \)收敛。
            \item 条件收敛的全部正项或全部负项级数一定发散，即若
            \item \(
                \sum_{n = 1}^{\infty} u_n
            \)条件收敛,则级数\(
                \sum_{n = 1}^{\infty} \frac{1}{2} (|u_n| + u_n)
            \)与级数\(
                \sum_{n = 1}^{\infty} \frac{1}{2} (u_n - |u_n|)
            \)
            发散。
        \end{enumerate}
    \end{conclusion}
    \begin{conclusion}{判别任意项级数收敛的步骤}{StepsToDiscriminateConvergenceOfSeries}
        \begin{enumerate}
            \item 判断\(\sum_{n = 1}^{\infty} |u_n|\)是否收敛。
            \item 若\(\sum_{n = 1}^{\infty} |u_n|\)收敛,则判断\(\sum_{n = 1}^{\infty} u_n\)的收敛性
            \item 若\(\sum_{n = 1}^{\infty} |u_n|\)发散,则判断\(\sum_{n = 1}^{\infty} u_n\)一定发散
        \end{enumerate}
        \boxed{\text{注}}莱布尼茨判别法只是判断交错级数收敛的一个充分条件，当莱布尼茨判别法失效时，应转回到正常的判别步骤 \\
        \boxed{\text{注}}绝对收敛级数与条件收敛级数的和必定是条件收敛级数
    \end{conclusion}
\subsection{幂级数}
\begin{mydef}{收敛点与收敛域}{ConvergencePointAndConvergenceDomain}
   对于函数项级数\(
       \sum_{n = 1}^{\infty} u_n(x)
    \)若在\(x_0 \in  I\)上\(
        \sum_{n = 1}^{\infty} u_n(x_0)
    \)收敛,则称\(x_0\)为\(  
        \sum_{n = 1}^{\infty} u_n(x)
    \)的收敛点,称\(I\)为\(
        \sum_{n = 1}^{\infty} u_n(x)
    \)的收敛域。称\(x_0 \notin I\)为\(
        \sum_{n = 1}^{\infty} u_n(x)
    \)的发散点。\(I'\)为\(
        \sum_{n = 1}^{\infty} u_n(x)
    \)的发散域。
\end{mydef}
\begin{mydef}{和函数}{SumFunction}
    设\(
        \sum_{n = 1}^{\infty} u_n(x)
    \)为函数项级数,若\(
        S(x) = \sum_{n = 1}^{\infty} u_n(x)
       (x \in  I)
    \),则称\(S(x)\)为\(
        \sum_{n = 1}^{\infty} u_n(x)
    \)的和函数。
\end{mydef}

\begin{mydef}{幂级数}{PowerSeries}
    形如\(
        \sum_{n = 1}^{\infty} a_n (x - x_0)^n
    \)的级数称为幂级数。其中常数\(a_n\)为系数，当\(x_0 = 0\)时，
    则\(
        \sum_{n = 1}^{\infty} a_n x^n
    \)为\(
        x
    \)的幂级数
\end{mydef}
\begin{conclusion}{幂级数的特点}{CharacteristicsOfPowerSeries}
    \begin{conclusion}{阿贝尔定理}{AbelTheorem}
        如果幂级数\(
            \sum_{n = 1}^{\infty} a_n (x - x_0)^n
        \)在\(x = x_0 \neq 0\)处绝对收敛，在\(
            x = x_1
        \)时发散
        则该幂级数在\(x < |x_0|\)上收敛，在\(
            x > |x_1|
        \)上发散
    \end{conclusion}
    \begin{mydef}{收敛区间与收敛半径}{ConvergenceIntervalAndConvergenceRadius}
        若\(
            \sum_{n = 1}^{\infty} a_n (x - x_0)^n
        \)在\(x \in (-R, R)\)上收敛，在\(x \in  (-\infty, R)\cup (R, \infty)\)上发散,
        则称\(R\)为收敛半径，\(I = (-R, R)\)为收敛区间。收敛点\(x = R\)或\(x = -R\)与收敛区间组成收敛域，
        如果只在\(x = 0\)处收敛，则收敛半径\(R = 0\),如果在整个数轴上收敛，则收敛半径为
        \(R = + \infty\)
    \end{mydef}
    如果\(R\)为幂级数的收敛半径，则在\(x \in (-R, R)\)上幂级数绝对收敛
\end{conclusion}

\begin{conclusion}{幂级数收敛半径与收敛域的求法}{ConvergenceRadiusAndConvergenceDomainOfPowerSeries}
    \begin{enumerate}
        \item 方法1,先求收敛半径得到收敛区间，再考察端点的敛散性
        \begin{enumerate}
            \item 求收敛半径
            \[
                \lim\limits_{n \to \infty} |\frac{a_{n+1}}{a_n}|(\lim\limits_{n \to \infty} \sqrt[n]{|a_n|}) = l
            \]
            \[
                R = \begin{cases}
                    \frac{1}{l} & 0 < l < +\infty \\
                    0 & l = +\infty \\
                    +\infty & l = 0 \\
                \end{cases}
            \]
            \item 讨论收敛区间端点处敛散性,即当\(0 < R < +\infty\)时,
            讨论\(sum_{n = 1}^{\infty} a_n (x - x_0)^n\)在\(x = -R\)和\(x = R\)处的敛散性
            \item 确定幂级数的收敛域
        \end{enumerate}
        \item 方法2,直接求收敛域
        求幂级数\(
            \sum_{n = 1}^{\infty} a_n (x - x_0)^n
        \)的收敛域,可将\(t = x - x_0\)代入幂级数,得到
        \[
            \sum_{n = 1}^{\infty} a_n t^n
        \]
        的收敛域与收敛半径
        \item 幂级数逐项求导或积分,收敛半径不变，但收敛域可能缩小或扩大
    \end{enumerate}
\end{conclusion}

\begin{conclusion}{幂级数的运算与和函数的性质}{PropertiesOfPowerSeries}
    \begin{enumerate}
        \item 设\(
            \sum_{n = 1}^{\infty} a_n x_n
        \)和\(
            \sum_{n = 1}^{\infty} b_n x_n
        \)为两个幂级数，收敛半径分别为\(R_1\)和\(R_2\)，则可以进行如下计算
        \begin{enumerate}
            \item \(\lambda \sum_{n = 1}^{\infty} a_n x_n + 
            \mu \sum_{n = 1}^{\infty} b_n x_n =
            \sum_{n = 1}^{\infty} (\lambda a_n + \mu b_n) x_n
        \)
            其中\(\lambda\)和\(\mu\)为任意常数
            \item \(
                \sum_{n = 0}^{\infty} a_n x_n \cdot \sum_{n = 0}^{\infty} b_n x_n =
                \sum_{n = 0}^{\infty} (c_n) x^n
            \),其中\(c_n = \sum_{k = 0}^{n} a_k b_{n - k}\),收敛半径\(
            R = min\{R_1, R_2\}(R_1 \neq R_2)\text{若}R_1 = R_2,\text{此结论不成立。}
            \)
        \end{enumerate}
        \item 设\(
            \sum_{n = 1}^{\infty} a_n x_n
        \)的收敛半径为\(R(R > 0)\),收敛区间为\(I = (-R, R)\),收敛域为\(D\)，
        则\(
            \sum_{n = 1}^{\infty} a_n x_n
        \)的和函数为
        \(
            S(x) = \sum_{n = 1}^{\infty} a_n x_n
        \),有如下性质成立
        \begin{enumerate}
            \item 和函数\(S(x)\)在\(I\)内任意次求导，并有逐项求导公式
            \[
                S^{k}(x) = \sum_{n = 1}^{\infty} a_n n^k x^{n - k}
            \]
            它的收敛半径仍为\(R\) 
            \item 和函数\(S(x)\)在\(I\)内任意次积分，并有逐项积分公式
            \[
                S(x) = \sum_{n = 1}^{\infty} a_n \frac{x^{n + 1}}{n + 1}
            \]
            它的收敛半径仍为\(R\)
            \item 若\(
                S(x) = \sum_{n = 1}^{\infty} a_n x^n
            \)在\(x = R(-R)\)处收敛，则该幂级数
            \begin{enumerate}
                \item \(\lim\limits_{x \to R-0} S(x) = \sum_{n = 0}^{\infty} a_nR^n\) (
                    \(
                        \lim\limits_{x \to R-0} S(x) = \sum_{n = 0}^{\infty} a_nR^n
                    \)
                )
                \item 可以在\([0, R]\)(\([-R, 0]\))上逐项积分，即
                \[
                    \int_{0}^{R} S(x) dx = \sum_{n = 0}^{\infty} \frac{a_n}{n+1}R^{n + 1}
                    (
                        \int_{-R}^{0} S(x) dx = \sum_{n = 0}^{\infty} \frac{-a_n}{n+1}(-R)^{n + 1}
                    )
                \]
                \item 逐项求导后的级数\(
                    \sum_{n = 0}^{\infty} n a_n x^{n - 1}
                \)可能在\(x = R(-R)\)发散
            \end{enumerate}
            \item 若\(
                S(x) = \sum_{n = 0}^{\infty} a_n x^n
            \)在\(x = R(-R)\)处发散，则该幂级数逐项求导后的级数
            \[
                \sum_{n = 0}^{\infty} n a_n x^{n - 1}
            \]
            也在\(x = R(-R)\)处发散， \\
            但逐项积分后的级数 
            \[
                \sum_{n = 0}^{\infty} \frac{a_n}{n+1} x^{n + 1}
            \]
            在\(x = R(-R)\)处有可能收敛
        \end{enumerate}
    \end{enumerate}
\end{conclusion}
\begin{mydef}{幂级数求和与幂级数展开式}{PowerSeriesSumAndExpansion}
    设\(f(x) = \sum_{n = 0}^{\infty} a_n(x = x_0)^n\),\(|x - x_0| < R\)，
    已知左端，求右端，是幂级数求和；已知右端，求左端，是幂级数展开式。
\end{mydef}

\begin{conclusion}{泰勒级数与麦克劳林级数}{TaylorAndMaclaurinSeries}
    设函数\(f(x)\)在点\(x_0\)的某一邻域内具有任意阶导数，则称级数
    \[
        \begin{aligned}
        &\sum_{n = 0}^{\infty} \frac{f^{(n)}(x_0)}{n!} (x - x_0)^n \\
        &= f(x_0) + \frac{f'(x_0)}{1!} (x - x_0) + \frac{f''(x_0)}{2!} (x - x_0)^2 + \cdots 
        + \frac{f^{(n)}(x_0)}{n!} (x - x_0)^n + \cdots
        \end{aligned}
    \]
    为\(f(x)\)的泰勒级数，\(x_0\)为展开点。
    若\(x_0 = 0\)，则称该级数为麦克劳林级数。
    \[
        \sum_{n = 0}^{\infty} \frac{f^{(n)}(0)}{n!} x^n = 
        f(0) + \frac{f'(0)}{1!} x + \frac{f''(0)}{2!} x^2 + \cdots
        + \frac{f^{(n)}(0)}{n!} x^n + \cdots
    \]
    展开后的级数是否收敛还需要验证
\end{conclusion}

\begin{conclusion}{函数展成幂级数的条件}{requirement for unfold}
    函数\(f(x)\)展开成幂级数：
    \[
        \begin{aligned}
        &f(x) = \sum_{n = 0}^{\infty} a_n (x - x_0)^n (|x - x_0| < R) \\
        &\Leftrightarrow \\
        &\text{在}(-R, R)\text{级数收敛于}f(x) \\
        &\Leftrightarrow \\
        &\lim\limits_{n \to \infty} R_n(x) = 0, \\
        &R_n(x) =  \frac{f^{n+1}[x_0 + \theta(x - x_0)]}{(n+1)!}(x - x_0)^{n+1},(0 < \theta < 1)\\
        &\text{是}f(x)\text{的拉格朗日余项}
        \end{aligned}
    \]
    \(f(x)\)在点\(x_0\)的幂级数展开式是唯一的，即若
    \(
        f(x) = \sum_{n = 0}^{\infty} a_n (x - x_0)^n (|x - x_0| < R)
    \)
    则
    \(
        a_n = \frac{f^{(n)}(x_0)}{n!}
    \)
    \\
    若存在常数\(L,M\)，使得\(\forall n \in \mathbb{N}\)，
    \(
        |f^{(n)}(x)| \leq LM^n, |x - x_0| \leq R
    \)
    则在\(x - x_0 < R\)上泰勒级数收敛于\(f(x)\),且
    \(
        \lim\limits_{n \to \infty} R_n(x) = 0
    \)
\end{conclusion}

\begin{formula}{常见函数的麦克劳林展开式}{usual marcliaurin unfold}
    \begin{enumerate}
        \item \[
            e^x = \sum_{n = 0}^{\infty} \frac{x^n}{n!} = 1 + x + \frac{x^2}{2!} + \cdots + \frac{x^n}{n!} + \cdots,(-\infty < x < +\infty)
        \]
        \item \[
            \sin x = \sum_{n = 0}^{\infty} \frac{(-1)^n x^{2n + 1}}{(2n + 1)!} = x - \frac{x^3}{3!} + \frac{x^5}{5!} - \frac{x^7}{7!} + \cdots,(-\infty < x < +\infty)
        \]
        \item \[
            \cos x = \sum_{n = 0}^{\infty} \frac{(-1)^n x^{2n}}{(2n)!} = 1 - \frac{x^2}{2!} + \frac{x^4}{4!} - \frac{x^6}{6!} + \cdots,(-\infty < x < +\infty)
        \]
        \item \[
            \ln(1 + x) = \sum_{n = 1}^{\infty} \frac{(-1)^{n + 1} x^n}{n},(-1 < x \leq 1)
        \]
        \item \[
            (1+x)^{\alpha} = 1 + \alpha x + \frac{\alpha(\alpha - 1)}{2!} x^2 + \cdots + \frac{\alpha(\alpha -1) \cdots
                (\alpha - n + 1)
            }{n!} x^n + \cdots
        \]
        ,\((-1 < x < 1)\)
    \end{enumerate}
    \((1+x)^{\alpha}\)的级数在\(x = \pm 1\)处的收敛性依\(\alpha\)的不同而异，特别，当\(\alpha = -1\)时，
    \[
        \begin{aligned}
            &\frac{1}{1+x} = \sum_{n = 0}^{\infty} (-1)^{n}x^{n} = 
            1 - x + x^2 - x^3 + \cdots + (-1)^{n}x^{n},(-1 < x <1)\\
            &\frac{1}{1-x} = \sum_{n = 0}^{\infty} x^{n} = 
            1 + x + x^2 + x^3 - \cdots + x^{n},(-1 < x <1)\\
        \end{aligned}
    \]
\end{formula}

\begin{conclusion}{幂级数求和与幂级数展开式的方法}{method of sum and unfold}
    \begin{enumerate}
        \item 直接法
        直接计算泰勒级数，证明\(\lim\limits_{n \to \infty} = 0\),即可求得它的部分和的极限
        \item 间接法
        利用上文常见的麦克劳林级数展开式，通过变量替换、或逐项求导，积分，待定性质，以及利用级数性质的方法
        来得到幂级数的和函数与展开式的方法. \\
        在求导或积分过程中，收敛半径不变，但端点处的敛散性可能会发生变化 \\
    \end{enumerate}
\end{conclusion}
\subsection{傅里叶级数}
\begin{mydef}{三角函数系的正交性}{triangle zhengjiao}
    三角函数系
    \[
        \{
            1, \cos \frac{\pi}{l} x , \sin \frac{\pi}{l} x ,
            \cos \frac{2\pi}{l} x , \sin \frac{2\pi}{l} x ,
            \cdots,
            \cos \frac{n\pi}{l} x , \sin \frac{n\pi}{l} x ,
            \cdots
        \}
    \]
    在区间\([-l, l]\)上是正交的，即
    \[
        \int_{-l}^{l} f(x) g(x) dx  = 0; f(x),g(x) \in \text{三角函数系}
    \]
    三角函数系中任何两个不同函数的乘积在\([-l, l]\)上的积分为0，即
    \[
        \begin{aligned}
            &\int_{-l}^{l} \cos \frac{n\pi}{l}x\mathrm{d}x = 
        \int_{-l}^{l} \sin \frac{n\pi}{l}x\mathrm{d}x = 0, \\
            &\int_{-l}^{l} \cos \frac{n\pi}{l}x \sin \frac{m\pi}{l}x \mathrm{d}x = 0 \\
            &\int_{-l}^{l} \cos \frac{m\pi}{l}x \cos \frac{n \pi}{l}x \mathrm{d}x
            =
            \int_{-l}^{l} \sin \frac{m\pi}{l}x \sin \frac{n \pi}{l}x \mathrm{d}x = 0
        \end{aligned}
    \]
\end{mydef}
\begin{mydef}{傅里叶系数与傅里叶级数}{Fourie series}
    若函数\(f(x)\)在\([-l, l]\)上可积分，且只定义在\([-l, l]\)上，或以\(2l\)为周期，
    则\(f(x)\)可以表示为
    \[
        f(x) = \frac{a_0}{2} + \sum_{n = 1}^{\infty} (a_n \cos \frac{n\pi}{l}x + b_n \sin \frac{n\pi}{l}x)
    \]
    称为\(f(x)\)的傅里叶级数.记作
    \[
        f(x) \sim \frac{a_0}{2} + \sum_{n = 1}^{\infty} (a_n \cos \frac{n\pi}{l}x + b_n \sin \frac{n\pi}{l}x)
    \]
    若
    \[
        \begin{cases}
            a_n = \frac{1}{l} \int_{-l}^{l} f(x) \cos \frac{n\pi}{l}x \mathrm{d}x, \\
            b_n = \frac{1}{l} \int_{-l}^{l} f(x) \sin \frac{n\pi}{l}x \mathrm{d}x.
        \end{cases}
    \]
    \(a_n, b_n\)称为\(f(x)\)的傅里叶系数.
\end{mydef}
\begin{conclusion}{傅里叶级数的收敛性(狄利克雷收敛定理(充分条件))}{dirichlet def}
    若\(f(x)\)在区间\(l\)上满足
    \begin{enumerate}
        \item 连续，或只有有限个第一类间断点 
        \item 只有有限个极值点
    \end{enumerate}
    则\(f(x)\)在\([-l, l]\)上的傅里叶级数收敛，且
    \[
        \begin{aligned}
            &\frac{a_0}{2} + \sum_{n = 1}^{\infty}(a_n \cos \frac{n\pi}{l}x + b_n \sin \frac{n\pi}{l}x) \\
            & \begin{cases}
            f(x), \text{若}x \in (-l, l)\text{为}f(x)\text{的连续点},  \\
            \frac{1}{2}\left[f(x+0) + f(x-0)\right], \text{若}x \in (-l, l)\text{且为}x\text{的第一类间断点} \\
            \frac{1}{2}\left[f(-l+0) + f(l-0)\right], \text{若}x = \pm l
        \end{cases}
        \end{aligned} 
    \]
\end{conclusion}

\begin{conclusion}{周期与非周期函数的傅里叶级数}{period or unperiod fourier}
    对于非周期函数,可以进行延拓，使得称为傅里叶级数
    \begin{enumerate}
        \item \(f(x)\)只在\([-l, l]\)上定义,进行周期延拓
        \item 只在\([0,l]\)有定义的函数,进行奇延拓或偶延拓,再作周期延拓
        \begin{enumerate}
            \item 奇函数 \\
            \(f(-x) = f(x)\),其傅里叶级数为正弦级数，即\(a_n = 0, n = 0, 1,2,\cdots\)\\
            \[
                f(x) \sim \sum_{i - 1}^{\infty} b_n sin \frac{n\pi}{l}x,
                b_n = \frac{2}{l} \int_{0}^{l} f(x) \sin \frac{n\pi}{l}x \mathrm{d}x
            \]
            \item 偶函数 \\
            \(f(x) = f(-x)\),其傅里叶级数为余弦级数，即\(b_n = 0, n = 0, 1,2,\cdots\)\\
            \[
                f(x) \sim \sum_{i - 1}^{\infty} a_n cos \frac{n\pi}{l}x,
                a_n = \frac{2}{l} \int_{0}^{l} f(x) \cos \frac{n\pi}{l}x \mathrm{d}x
            \]
        \end{enumerate}
        \item 在进行延拓时,奇延拓出来的\(a_0 = 0\),偶延拓出来的\(a_0 \neq 0\)，所以求偶延拓时要注意加上\(\frac{a_0}{2}\)
    \end{enumerate}
\end{conclusion}

\begin{conclusion}{函数\(f(x)\)的傅里叶展开式}{fx unfold fourier}
    设\(f(x)\)在\([-l, l]\)上，并满足狄利克雷条件，
    \begin{enumerate}
        \item  若有\(f(x)\)在\((-l, l)\)上连续，则\(f(x)\)可展成傅里叶级数，即
        \[
        f(x) = \frac{a_0}{2} + \sum_{n = 1}^{\infty} (a_n \cos \frac{n\pi}{l}x + b_n \sin \frac{n\pi}{l}x),
        (-l < x < l)
        \]
        \item 若\(f(x)\)在\([-l, l]\)上连续且\(f(-l) = f(l)\),则\(f(x)\)在\([-l, l]\)上可展成
        傅里叶级数,即上式在\([-l, l]\)上成立
    \end{enumerate}
\end{conclusion}

\subsection{普适方法}

\begin{conclusion}{常数项级数敛散性的判定}{const term series}
    \begin{enumerate}
        \item 判别级数\(\sum_{n = 1}^{\infty} a_n\)的类型
        \item 若\(\sum_{n = 1}^{\infty} a_n\)是正项级数，则首先考察\(
        \lim\limits_{n \to \infty} a_n
        \begin{cases}
            \neq 0, \text{级数发散} \\
            = 0, \text{进一步判断}
        \end{cases}
        \)
        \begin{enumerate}
            \item 一般项含有\(n!\)或者\(n\),选择比值判别法
            \item 一般项含有以\(n\)为指数的因子，采用根值判别法
            \item 含有形如\(n^{\alpha}\)的因子，一般选择比较判别法
            \item 若一般项\(u_n\)可表示为\(f(n)\)，\(f(x)\)在\(
                [1, +\infty)
            \)单调下降,且其积分的敛散性容易判断，可使用积分判别法
            \item 利用已知敛散性结果，结合级数性质，进行判断
            \item 如上述方法均失效，采用定义，考察\(\lim\limits_{n \to \infty}S_n\)
            是否存在，即判断\(S_n\)是否有上界
        \end{enumerate}
        \item 若为任意项级数，先考察\(
        \lim\limits_{n \to \infty} a_n
        \begin{cases}
            \neq 0, \text{级数发散} \\
            = 0, \text{进一步判断}
        \end{cases}
        \)
        若\(\lim\limits_{n \to \infty} a_n = 0\),则如下进行
        \begin{enumerate}
            \item 按正项级数判别法，判断\(
                \sum_{n = 1}^{\infty} |a_n|
            \)是否收敛，若收敛，则\(\sum_{n = 1}^{\infty} a_n\)绝对收敛
            \item 若\(
                \sum_{n = 1}^{\infty} |a_n|
            \)发散，则看其是否为交错级数，若是，则采用莱布尼茨判别法，若收敛，则
            则\(\sum_{n = 1}^{\infty} a_n\)条件收敛
            \item 若\(\sum_{n = 1}^{\infty} a_n\)为交错级数，但不满足莱布尼茨判别法条件
            或\(\sum_{n = 1}^{\infty} a_n\)不是交错级数，则有如下方法
            \begin{enumerate}
                \item 比值判别法,当\(\lim\limits_{n \to \infty} \frac{|u_{n+1}|}{u_{n}} = \rho > 1 \)时，
                级数\(sum_{n = 1}^{\infty} u_n\)发散
                \item 根植判别法,当\(\lim\limits_{n \to \infty} \sqrt[n]{|u_n|} = \rho > 1 \)时，
                级数\(sum_{n = 1}^{\infty} u_n\)发散
                \item 讨论\(\{S_{2n}\},S_{2n+1}\)的敛散性，以及其他能够判定部分和数列\(\{
                    S_n
                \}\)敛散性的方法
            \end{enumerate}
        \end{enumerate}
             \item 结合级数性质判定敛散性
    \end{enumerate}
\end{conclusion}

\begin{conclusion}{常数项级数求和}{const term series sum}
    \begin{enumerate}
        \item 直接计算部分和\(S_n = \sum_{n = 1}^{\infty} a_n\)得到级数的和
        \(
            \sum_{n = 1}^{\infty} a_n = \lim\limits_{n \to \infty} S_n
        \)
        \item 利用分解法和已知的幂级数展开式
        \item 转化为求幂级数的和函数
        \item 借助于函数的傅里叶级数展开式
    \end{enumerate}
\end{conclusion}

\begin{conclusion}{求幂级数展开式与利用幂级数展开式求和}{power series}
    \begin{enumerate}
        \item 函数
        \item 求傅里叶级数展开式
        \item 根据傅里叶级数特性，确定要计算的常数项系数
        \item 根据收敛定理求常数项系数之和
    \end{enumerate}
\end{conclusion}

\begin{conclusion}{\(S(X)\)的导数表示}{S(X) derivative}
    \[
        S(x) = \int_{0}^{x} f(t) dt + S(0)
    \]
    \begin{center}
        \boxed{\text{注意不要忘记S(0)}}
    \end{center}
\end{conclusion}


